%% kapa: exam=1 round=28 year=2017 period=1 form=A subject=economics subject_ko=경제학원론 seq=22 number=62
%% source: Q-Net 감정평가사 기출문제 (cst00309) artlSeq=5199603 file=extracted/1008/5199603/2017년 제28회 감정평가사 1차시험 문제지/1교시 A형.pdf
%% conversion: kordoc 4.24.51 -> tools/gen1.py (auto); review=proofread-sonnet; flags=none
\begin{question}{62}
경기변동이론에 관한 설명으로 옳은 것은?
\begin{choices}
  \item 실물경기변동이론(real business cycle theory)은 통화량 변동 정책이 장기적으로 실질 국민소득에 영향을 준다고 주장한다.
  \item 실물경기변동이론은 단기에는 임금이 경직적이라고 전제한다.
  \item 가격의 비동조성(staggered pricing)이론은 새고전학파(New Classical) 경기변동이론에 포함된다.
  \item 새케인즈학파(New Keynesian) 경기변동이론은 기술충격과 같은 공급충격이 경기변동의 근본 원인이라고 주장한다.
  \item 실물경기변동이론에 따르면 불경기에도 가계는 기간별 소비선택의 최적조건에 따라 소비를 결정한다.
\end{choices}
\end{question}
\answer{5}
